\chapter{Bayesian Optimization}

\green{This chapter (and the rest of this work) borrows heavily from the book \cite{garnett_bayesoptbook_2023} as a source of consistent terminology.} This book might be the most approachable and complete treatment of Bayesian optimization the author has found. We rely on it for standardizing notation and for some of the proofs and concepts, as well as for a general sketch of how this chapter is set out.

\section{Introduction}

\green{Bayesian optimization is central to the experimental framework developed in this work. Understanding its motivation, assumptions, and mechanics is therefore essential before proceeding, and this section aims to provide a clear and self-contained foundation. The chapter that follows prioritizes conceptual clarity over generality and formal definitions with algorithmic details follow in subsequent chapters.}

\green{The necessity for Bayesian optimization arises from a particular kind of constraint, one that is distinct from the structural constraints we typically encounter in classical optimization. In traditional optimization, constraints govern the \textit{form} of the solution space: for instance, in integer programming, the objective function is defined only over discrete inputs, or we restrict our interest to solutions that are integer-valued. In these settings, the primary objective is to locate an optimum of some objective function $f$ over a well-defined domain $\mathcal{X}$. The computational cost of doing so, measured in the number of function evaluations required, is a secondary concern\footnote{Sometimes this is not even a concern. We evaluate these functions in terms of the amout of wall-clock time it takes to converge to a solution, not necessarily the amount of computations required to arrive at that solution.}.}

\green{Bayesian optimization shares this goal but introduces a critical additional assumption: the objective function $f$ is \textit{expensive to evaluate}. That is, each query $f(\mathbf{x})$ carries a significant cost: computational, financial, physical, or temporal. In this scenario, exhaustive or gradient-based search over $\mathcal{X}$ becomes impractical. Under this paradigm minimizing the number of evaluations is no longer a secondary concern but a core objective in its own right. We therefore seek not only to find the optimum of $f$, but to do so in as few evaluations as possible.}

\green{To make this concrete, consider the problem of hyperparameter optimization in machine learning, which will become directly relevant to this work in later chapters. When training a neural network, one must specify a learning rate $\eta$ which is the rate at which the model updates its weights in response to the gradient of the loss. If $\eta$ is set too high, training may become unstable and fail to converge. If it is set too low, convergence will be prohibitively slow or the model may stagnate at a poor local minimum. Evaluating a single candidate value of $\eta$ requires a full training run, which for modern large-scale models can take hours or days. Exhaustive grid search or random search over the hyperparameter space is therefore wasteful and sometimes impossible. A more principled approach is needed.}

\green{Bayesian optimization provides exactly this. Rather than treating $f$ as a black box to be sampled blindly, we maintain a probabilistic \textit{surrogate model} that encodes our current beliefs about the shape of $f$ given the evaluations observed so far. At each iteration, an \textit{acquisition function} uses this surrogate to identify the point in $\mathcal{X}$ that is most worth evaluating next, balancing the competing demands of \textit{exploitation} (querying near known good regions) and \textit{exploration} (querying in regions of high uncertainty). This loop composed of update the surrogate, optimize the acquisition function, evaluate $f$, and finally start again, continues until a convergence criterion is met or a budget of evaluations is exhausted.}

\green{There are many domains beyond hyperparameter tuning where this paradigm is natural and necessary: experimental design in the physical sciences, A/B testing in product development, materials discovery, and the calibration of expensive simulation models, among others. What these problems share is the same fundamental structure: an objective value or function that is cheap to specify but costly to evaluate, with no available gradient information and a need to reach a good solution within a limited query budget.}

\green{The remainder of this chapter develops the mathematical machinery required to make this framework precise. We begin with Gaussian processes as the surrogate model of choice, before turning to the construction and properties of standard acquisition functions, and finally to the full Bayesian optimization algorithm and its convergence behavior.
}
\green{As the models that we build and the problems we solve in general become more and more complicated, there has been a growing interest in this area of statistics as a solution to generating viable extrema candidates of expensive to evaluate functions. There has been extensive work on this field recently and in the past few decades, with examples such as \cite{bergstra2013hyperopt}, \cite{balandat2020botorch}, and \cite{akiba2019optuna} that have impulsed this filed considerably. Hyperparameter optimization has recently been talked about as the \say{killer app} for Bayesian optimization, which revitalized interest in Gaussian Processes in the context of machine learning and Bayesian optimization in general. The interest in the field has lead to a great amount of interest and therefore research and publications tackling the method and providing theoretical backing to it. }

\green{In this chapter we will show the basic theoretical structure of Bayesian optimization, the three central components that make it up, and the decisions we have to make to be able to approximate these functions from a Bayesian point of view. }

\section{Components of BO}

\green{Before describing the components of Bayesian optimization in detail, it is worth being precise about the class of functions we are concerned with. A function $f: \mathcal{X} \to \mathbb{R}$ is said to be a \textit{black-box function} if its internal structure is either unknown or computationally infallible to observe directly. More precisely, we have no access to an analytical expression for $f$, no gradient information $\nabla f(\mathbf{x})$, and no closed-form characterization of its global behavior. The only information available to us is the result of directly querying the function: given an input $\mathbf{x} \in \mathcal{X}$, we may observe the output $f(\mathbf{x})$. This stands in contrast to the functions typically encountered in classical optimization, where gradient-based methods can exploit the analytic structure of $f$ to navigate the search space efficiently. In the black-box setting, no such structure is available, and evaluations themselves are the sole source of information about $f$.}


Bayesian optimization is a sequential model-based approach to optimizing black-box functions that are expensive to evaluate. The goal is to find in some domain $\mathcal{X}$ a value $x$ that maximizes the objective function $f$ with the least possible number of evaluations. Practically, this means \green{evaluating} the value $f(x)$ for the least amount of times possible. One important distinction from other areas of optimization is that we do not impose any restrictions on the function $f$, except continuity\footnote{As we will see in later chapters, the entire point of this work will be relaxing this restriction.}. \green{This is notably different from many classical settings. In integer optimization for instance, the domain $\mathcal{X}$ is restricted to $\mathbb{Z}^d$, meaning the objective function is only defined (or only of interest) at discrete points, and the geometry of the search space changes fundamentally as a result. At the other extreme are methods such as the finite element method (FEM) used to solve partial differential equations numerically. They require not only that the functions involved be continuous, but that they belong to appropriate Sobolev spaces, possessing weak derivatives of sufficient order to guarantee the existence and stability of solutions.} Bayesian optimization asks for neither: $f$ need not be discrete, nor differentiable, nor analytically specified in any way except continuity. Formally, let $f$ be the objective function that we want to maximize, and let $\mathcal{D}$ be the set of observations we have made so far, which consists of pairs $(x_i, y_i)$ where $x_i$ is the input and $y_i$ is the corresponding output of the function $f(x_i)$. We want to find the input $x$ that maximizes the function $f(x)$, i.e., $x^* = \arg\max_\mathcal{X} f(x)$

The Bayesian optimization algorithm consists of three main components: a prior distribution over the objective function, an acquisition function, and a method for updating the prior distribution based on the observed data. The way this model works can be thought of as the following three steps: 

\begin{enumerate}
	\item Determining a probabilistic model that captures the observed behaviour of $f$, either by assuming a prior on the first step or going off data points $(x_i, f(x_i))$ for later iterations. We call this probabilistic mode of $f$ the prior, since it is used basically as one under this scheme\footnote{Another common way to jump-start this process is to begin by evaluating some random points in $\mathcal{X}$ and continuing the process from there. }.
	\item Optimizing the selected acquisition function that probabilistically estimates the best possible points for evaluation in the following iteration. 
	\item Evaluating the function at the best candidate for improvement, therefore getting new data and updating our probabilistic model of the behaviour of the function. 
\end{enumerate}

The prior distribution over the objective function $f(x)$ is usually assumed to be a Gaussian process, but that need not be the case. While these type of models are beyond the scope of this work, as examples, we refer to \cite{shah2013bayesian} for the prior set to a $t$-Student distribution and \cite{wangpoisson} for a case study of the Poisson prior. Given the set of observations $\mathcal{D}$, the selected prior distribution defines a posterior distribution that takes into account the uncertainty in the observations.

The acquisition function is a criterion that measures the expected utility of each point in the input space $x$ for the purpose of optimization. The most commonly used acquisition function is the expected improvement (EI), which measures the expected improvement in the objective function if we were to evaluate the function at a particular input $x$. Much like the case with the prior distributions, there are other choices such as Thompson sampling (TS) and Upper Confidence bound (UCF), references for each can be found in  \cite{kaufmann2012thompson} and \cite{kaufmann2012bayesian} respectively. 

\[ EI(x) = \mathbb{E}[max(f(x) - f(x_{best}), 0)], \]

where $x_{\text{best}}$ is the current best input found so far. The idea is to select the input $x$ that maximizes the EI, which balances exploration (trying new inputs) and exploitation (focusing on inputs that are likely to be good) this balance os controlled by hyperparameters, so that is another layer of complexity that needs to be controlled. In a sense, what we are doing is differing solving the main optimization problem by generating a simpler to solve optimization problem in finding the relevant extrema of the acquisition function, thus allowing us to spend as little function evaluations as possible on evaluating the target function $f$. 

The final component of BO is the method for updating the prior distribution based on the observed data. This is done under the overarching philosophy of Bayesian inference, which allows us to compute the posterior distribution over the objective function given the prior distribution and the observed data. The posterior distribution is then used as the prior distribution for the next iteration of the algorithm. What we wish to do is infer the value $\phi = f(x)$ for some unobserved point in the domain $\mathcal{X}$. 

\section{Formalization}

This process begins by setting the prior distribution $\mathbb{P}(\phi | x)$ which is our belief of possible values for $\phi$ before observation of the unknown $x$. 

Assuming we observe the objective function at $x$ and measure $y$, our optimization model supposes that the distribution of this measurement is determined by the value of interest $\phi$ through the observation model $\mathbb{P}(y | x, \phi)$. With this new value $y$ we can update our beliefs and generate the posterior function through the use of Bayes' theorem: 
\[ \mathbb{P}(\phi | x, y) = \frac{\mathbb{P}(\phi | x) \cdot \mathbb{P}(y | x, \phi)}{\mathbb{P}(y | x)}. \] 

As with all applications of Bayes' theorem such as this one, the denominator is simply a constant term used to scale the probability. This means that we can simplify the expression by only keeping the kernel of the distribution and dropping the normalizing term, and using the proportionality operator $\propto$, as follows. 
\[ \mathbb{P}(\phi | x, y) \propto \mathbb{P}(\phi | x) \cdot \mathbb{P}(y | x, \phi). \] 

In the Bayesian framework, the posterior distribution is often not the final goal but a starting point for further tasks such as prediction or decision-making. This builds on the iterative process under which we are working in Bayesian optimization, where new information is added to the approximation of the function every time a function evaluation happens. As an example, consider the case where we want to predict the outcome of an independent, repeated noisy observation at $x, y'$. We can treat the outcome as a random variable and derive its distribution by integrating our posterior belief about $\phi$ against the observation model:

\[ \mathbb{P}(y' | x, \mathcal{D}) = \int \mathbb{P}(\phi | x, \mathcal{D}) \cdot \mathbb{P}\left(y' | x, \phi\right) \, d \phi. \]

This equation represents the posterior predictive distribution for $y'$ given we know $x$ and a set of data $\mathcal{D}$. By integrating over all possible values of $\phi$, weighted by their plausibility, the posterior predictive distribution naturally accounts for uncertainty in the unknown objective function value.

From here is where the link to Gaussian Processes lies. Gaussian Processes serve as a natural bridge between the Bayesian framework and Bayesian optimization. They provide a flexible, non-parametric prior of the form $\mathbb{P}(f)$ over a certain function $f$, and their posterior predictive distributions can be efficiently computed and used to guide the search for the optimal input value. The acquisition function, which is derived from the GP's posterior predictive distribution, plays a crucial role in balancing exploration and exploitation during the optimization process.

\textcolor{red}{Acabo aqui o faltaria algo más?}
















